\section{Introduction}
\label{sec:intro}

Large Language Models (LLMs) generate text autoregressively: each new token requires a full forward pass conditioned on all preceding tokens, making inference latency proportional to the output length.
The resulting low GPU utilization and high user-perceived waiting time constitute a primary bottleneck in production LLM serving, particularly for latency-sensitive scenarios such as real-time conversational assistants and multi-turn agentic workflows.
Speculative decoding~\citep{pmlr-v202-leviathan23a,chen2023accelerating} offers a principled solution: a lightweight \textit{draft model} proposes a block of candidate tokens, and the full-size \textit{target model} verifies the entire block in a single forward pass via rejection sampling, accepting the longest prefix consistent with the target distribution and appending one bonus token.
Because verification is parallel and the acceptance rule preserves the target distribution exactly, speculative decoding accelerates generation without any quality loss.

The design of the draft model governs the trade-off between drafting latency and acceptance rate. Early drafters are autoregressive~\citep{li-etal-2024-eagle,cheng2024recurrentdrafterfastspeculative}, conditioning each position on previously sampled tokens. However, their drafting latency grows linearly with the block size, forcing these methods to use short blocks and shallow architectures. To break this sequential bottleneck, parallel drafters~\citep{pmlr-v235-cai24b,liu2026dartdiffusioninspiredspeculativedecoding,chen2026dflash} have emerged as a compelling alternative: all draft positions are produced in a single forward pass, making drafting latency nearly independent of block size. This structural advantage theoretically allows parallel drafters to efficiently generate substantially longer draft blocks.

However, fully unlocking the potential of large parallel draft blocks introduces two critical bottlenecks—one in generation quality, and the other in system efficiency.
First, because parallel drafters predict each position independently, they cannot model inter-token dependencies within a block. This independence leads to multi-modal collisions and rapid acceptance decay at later positions~\citep{gu2018nonautoregressive,pmlr-v162-huang22m}. 
Second, determining the optimal verification length remains a challenge. While parallel generation easily produces long draft blocks, indiscriminately verifying all proposed tokens degrades system throughput, particularly under high-concurrency workloads~\citep{liu2024turbospec,hu2026echo}. The ideal verification length varies along two axes. 
On the data side, structured requests like code naturally sustain higher acceptance rates than open-ended chat~\citep{xia-etal-2024-unlocking,abramovich2026speed}. 
On the system side, verifying extra tokens is nearly free under light loads. Under heavy loads, however, verifying tokens with a high rejection risk occupies critical batch capacity that could otherwise serve other active requests~\citep{liu2024optimizing, wu-etal-2025-tetris}.

To address these bottlenecks, we introduce \textbf{DSpark}, a speculative decoding framework that unifies high-throughput parallel generation with adaptive, load-aware verification. At its core, DSpark is designed to resolve the inherent trade-offs in draft generation and verification through two complementary mechanisms. 
\begin{itemize}
    \setlength{\itemsep}{10pt}
    \item First, to overcome the lack of inter-token dependencies, DSpark adopts a semi-autoregressive architecture. It keeps the computationally expensive draft backbone fully parallel, appending only a lightweight serial output head to inject local transition information. This design preserves the drafting speed of parallel models while significantly mitigating suffix decay.
    \item Second, to resolve the system-level bottleneck, DSpark employs confidence-scheduled verification. By coupling a confidence head—which estimates per-position prefix survival probabilities—with a hardware-aware scheduler, DSpark dynamically tailors the verification length for each request. This scheduler leverages real-time engine throughput profiles to route target verification budget only toward tokens with the highest expected return.
\end{itemize}

We extensively evaluate DSpark across both controlled offline benchmarks and production-scale online deployments.
On controlled offline benchmarks—spanning mathematical reasoning, code generation, and daily chat—DSpark consistently outperforms strong baselines. Specifically, across the Qwen3-4B, 8B, and 14B target models~\citep{yang2025qwen3}, it improves the macro-average accepted length over the autoregressive Eagle3~\citep{li2026eagle} by 30.9\%, 26.7\%, and 30.0\%, and over the parallel DFlash~\citep{chen2026dflash} by 16.3\%, 18.4\%, and 18.3\%, respectively. Beyond top-line metrics, our fine-grained position-wise analysis reveals the distinct generation characteristics of different drafters, empirically demonstrating how DSpark successfully combines the high initial-token capacity of parallel models with the suffix coherence of autoregressive models.

Beyond offline evaluation, we deployed DSpark within the DeepSeek-V4~\citep{deepseekai2026deepseekv4highlyefficientmilliontoken} serving system to assess its performance under live user traffic. Compared to the prior MTP-1 production baseline~\citep{liu2024deepseek}, DSpark significantly broadens the system's operational envelope. Specifically, it consistently accelerates per-user generation speeds by 60\%--85\% (V4-Flash) and 57\%--78\% (V4-Pro) at matched aggregate throughput capacities. 
Furthermore, under strict Service Level Agreements (SLAs) where the baseline's capacity deteriorates severely—such as 120 TPS for Flash and 50 TPS for Pro—DSpark mitigates verification overhead to maintain robust throughput. By overcoming this performance cliff, DSpark unlocks strict interactivity tiers that were previously unattainable, effectively shifting the Pareto frontier of LLM serving.

To foster collective advancement within the open-source community, we are making our artifacts publicly available. Specifically, we release the trained \href{https://huggingface.co/deepseek-ai/DeepSeek-V4-Pro-DSpark/tree/main}{DSpark checkpoints} for both the DeepSeek-V4-Flash~(preview) and DeepSeek-V4-Pro~(preview) models. Furthermore, we open-source \href{https://github.com/deepseek-ai/DeepSpec}{DeepSpec}, an algorithm-driven training repository, including Eagle3, DFlash and DSpark. 
These artifacts are intended to support further research on efficient LLM serving.